\chapter{Geriatric Medicine and Aging-Related Disorders}
\label{ch:geriatric}
Geriatric medicine addresses the unique health needs of older adults, a population that is growing rapidly worldwide. By 2050, the global population aged 60 and older is expected to exceed 2.1 billion. Aging is a complex biological process characterized by progressive loss of physiological integrity, leading to impaired function and increased vulnerability to disease. This chapter covers the biology of aging, geriatric syndromes, falls and mobility, cognitive disorders, polypharmacy, incontinence, sensory loss, osteoporosis, cardiovascular aging, and end-of-life care---all through a geriatric lens emphasizing function, independence, and quality of life.

%-------------------------------------------------------------
\section{Biology of Aging}
%-------------------------------------------------------------

%-------------------------------------------------------------

%-------------------------------------------------------------

The biology of aging encompasses progressive, intrinsic physiological changes occurring across organ systems over time. Processes include cellular senescence, genomic instability, telomere attrition, impaired proteostasis, and mitochondrial dysfunction that diminish physiological reserve. These biological hallmarks increase susceptibility to frailty, multimorbidity, altered pharmacokinetics, and functional decline in older adults.

%-------------------------------------------------------------

%-------------------------------------------------------------

%-------------------------------------------------------------

%-------------------------------------------------------------

\label{sec:Biology of Aging}
%-------------------------------------------------------------%-------------------------------------------------------------

\subsection{Cellular Senescence}
\label{sec:Cellular Senescence}

Cellular senescence is a state of irreversible cell cycle arrest in response to various stressors including telomere shortening, DNA damage, oncogene activation, and oxidative stress. Senescent cells accumulate in multiple tissues with age and secrete a complex mixture of pro-inflammatory cytokines, chemokines, growth factors, and matrix metalloproteinases---the senescence-associated secretory phenotype (SASP). The SASP drives chronic inflammation, disrupts tissue architecture, and contributes to age-related diseases. Key pathways include p53/p21\verb|^|CIP1 and p16\verb|^|INK4a/Rb. Mouse studies demonstrate that selective elimination of senescent cells (using senolytic agents such as dasatinib plus quercetin, or fisetin) attenuates age-related pathology and extends healthspan. Senolytics are being actively investigated in human clinical trials for frailty, osteoarthritis, and other conditions.

\subsection{Frailty as a Syndrome}
\label{sec:Frailty as a Syndrome}

Frailty is a clinically recognizable state of increased vulnerability to stressors resulting from age-related decline in physiological reserve across multiple organ systems. It predicts adverse outcomes including falls, hospitalization, institutionalization, and death, independent of chronological age and comorbidity. Frailty is a dynamic process that can improve or worsen over time, making early detection and intervention clinically important.

\subsection{Theories of Aging}
\label{sec:Theories of Aging}

Numerous theories attempt to explain the molecular and cellular basis of aging. The \textbf{telomere shortening theory} posits that progressive shortening of telomeres---the protective caps at chromosome ends---with each cell division eventually triggers replicative senescence or apoptosis. Telomerase, the enzyme that elongates telomeres, is silenced in most somatic cells. The \textbf{free radical theory} (Harman, 1956) proposes that reactive oxygen species (ROS) generated by mitochondrial metabolism cause cumulative oxidative damage to DNA, proteins, and lipids, driving cellular aging. The \textbf{mitochondrial theory} extends this concept, focusing on the accumulation of mitochondrial DNA mutations and impaired oxidative phosphorylation with age. The \textbf{epigenetic theory} emphasizes age-related changes in DNA methylation, histone modification, and chromatin remodeling that alter gene expression patterns without changing the DNA sequence. Epigenetic clocks (e.g., Horvath clock) use DNA methylation patterns to estimate biological age. \textbf{Inflammaging} describes the chronic, low-grade, sterile inflammatory state that characterizes aging, driven by activation of the innate immune system, accumulation of damage-associated molecular patterns (DAMPs), and dysregulation of cytokine networks (elevated IL-6, TNF-alpha, CRP). This promotes many age-related diseases including atherosclerosis, diabetes, dementia, and frailty.

%-------------------------------------------------------------
\section{Cardiovascular Aging}
%-------------------------------------------------------------

%-------------------------------------------------------------

%-------------------------------------------------------------

%-------------------------------------------------------------

%-------------------------------------------------------------

%-------------------------------------------------------------

%-------------------------------------------------------------

\label{sec:Cardiovascular Aging}
%-------------------------------------------------------------%-------------------------------------------------------------

\subsection{Aortic Stenosis}
\label{sec:geriatric:Aortic Stenosis}
Calcific (degenerative) aortic stenosis (AS) is the most common valvular heart disease in older adults, affecting 3--5\% of those over 75. It is characterized by progressive thickening, fibrosis, and calcification of the aortic valve cusps, driven by lipid deposition, inflammation, and osteogenic differentiation of valvular interstitial cells. Once symptomatic (classic triad: dyspnea, angina, syncope), prognosis without intervention is poor (50\% mortality at 2 years). In older adults, transcatheter aortic valve replacement (TAVR) has largely replaced surgical aortic valve replacement and is the standard of care for most elderly patients, with excellent outcomes and shorter recovery times.

\subsection{Atrial Fibrillation in the Elderly}
\label{sec:Atrial Fibrillation in the Elderly}
Atrial fibrillation (AF) prevalence increases from 2\% in adults under 65 to 15--20\% in those over 80. Geriatric-specific management challenges include higher thromboembolic risk, increased bleeding risk, renal impairment affecting DOAC dosing, high rates of polypharmacy, and increased risk of falls complicating anticoagulation decisions. Anticoagulation for stroke prevention is indicated for most older adults with AF, using DOACs in preference to warfarin due to lower intracranial hemorrhage risk. Rate control is the default strategy in older adults, reserving rhythm control for symptomatic patients despite adequate rate control.

\subsection{Heart Failure with Preserved Ejection Fraction}
\label{sec:geriatric:Heart Failure with Preserved Ejection Fraction}
Heart failure with preserved ejection fraction (HFpEF, EF $\ge$50\%) is the predominant heart failure phenotype in older adults, especially older women with hypertension, diabetes, obesity, and atrial fibrillation. Pathophysiology includes left ventricular diastolic dysfunction, systemic inflammation, pulmonary hypertension, and chronotropic incompetence. Clinical presentation is dyspnea on exertion, exercise intolerance, and fluid overload. Diagnosis requires clinical signs/symptoms of HF, preserved LVEF, and evidence of diastolic dysfunction or elevated filling pressures. Management includes diuretics for volume overload, blood pressure control, SGLT2 inhibitors, and management of comorbidities. Prognosis is variable.

\subsection{Hypertension in the Elderly}
\label{sec:Hypertension in the Elderly}

Hypertension affects over 70\% of adults aged 65 and older. The predominant form in older adults is isolated systolic hypertension (ISH), defined as systolic BP $\ge$130 mmHg with diastolic BP <80 mmHg, resulting from age-related arterial stiffening. ISH confers significant risk for stroke, coronary artery disease, heart failure, and cognitive decline. The SPRINT trial demonstrated that intensive systolic BP lowering to <120 mmHg reduced cardiovascular events and mortality in adults over 75, though with increased risk of orthostatic hypotension, electrolyte abnormalities, and acute kidney injury. Current guidelines recommend a target of <130/80 mmHg for most older adults, but treatment decisions should be individualized based on frailty, falls risk, polypharmacy, and life expectancy. First-line agents include thiazide diuretics, calcium channel blockers, and ACE inhibitors/ARBs. Initial doses should be low with gradual up-titration.

%-------------------------------------------------------------
\section{Cognitive and Neurologic Disorders of Aging}
%-------------------------------------------------------------

%-------------------------------------------------------------

%-------------------------------------------------------------

%-------------------------------------------------------------

%-------------------------------------------------------------

%-------------------------------------------------------------

%-------------------------------------------------------------

\label{sec:Cognitive and Neurologic Disorders of Aging}
%-------------------------------------------------------------%-------------------------------------------------------------

\subsection{Alzheimer Disease}
\label{sec:Alzheimer Disease}
Alzheimer disease (AD) is a progressive neurodegenerative disorder and the most common cause of dementia (60--80\% of cases), with prevalence doubling every 5 years after age 65. In the geriatric population, AD presents with a higher burden of comorbidity, polypharmacy, and functional dependence. Late-onset AD (>65 years) is distinguished from early-onset AD by its sporadic nature, strong association with the APOE epsilon-4 allele, and slower progression. Clinical features include progressive memory loss, impaired executive function, language difficulties, and behavioral and psychological symptoms of dementia (BPSD: agitation, aggression, psychosis, depression, apathy). Diagnosis is clinical based on history, examination, and cognitive testing. Management includes non-pharmacological approaches first (person-centered care, environmental modification, structured routines); antipsychotics carry a black box warning for increased mortality in elderly dementia patients and should be used only for severe, refractory symptoms. Prognosis is progressive and ultimately fatal; end-of-life considerations include advance care planning and hospice referral.

\subsection{Delirium}
\label{sec:Delirium}

Delirium is an acute, fluctuating disturbance in attention, awareness, and cognition, caused by an underlying medical condition, intoxication, or medication side effect. It is a medical emergency affecting 10--30\% of hospitalized older adults and up to 50\% of postoperative patients. Delirium is associated with increased mortality, longer hospital stays, functional decline, and incident dementia. Three subtypes exist: hyperactive (agitation, restlessness), hypoactive (withdrawal, lethargy, reduced speech---often missed or misdiagnosed as depression), and mixed. The Confusion Assessment Method (CAM) is the most widely used diagnostic tool: acute onset and fluctuating course, plus inattention, plus either disorganized thinking or altered level of consciousness. Predisposing factors include advanced age, dementia, sensory impairment, polypharmacy, and multiple comorbidities. Precipitating factors include infection, electrolyte disturbance, surgery, psychoactive medications, pain, sleep deprivation, and urinary retention. Prevention is paramount (the Hospital Elder Life Program---HELP). Pharmacologic treatment (haloperidol or atypical antipsychotics) is reserved for severe agitation that threatens safety. Prognosis is variable; delirium is associated with increased mortality and functional decline.

\subsection{Dementia with Lewy Bodies}
\label{sec:Dementia with Lewy Bodies}
Dementia with Lewy bodies (DLB) is a progressive neurodegenerative disorder accounting for 5--15\% of dementia cases. Core clinical features include fluctuating cognition with pronounced variations in attention and alertness, recurrent well-formed visual hallucinations (often detailed, involving people or animals), and parkinsonism (bradykinesia, rigidity, gait disorder). REM sleep behavior disorder (RBD) is a supportive feature, often predating dementia onset by years. In the geriatric population, DLB patients are extremely sensitive to neuroleptics (neuroleptic sensitivity reaction---severe parkinsonism, sedation, autonomic instability, increased mortality), which should be avoided. Management includes cholinesterase inhibitors (rivastigmine is FDA-approved for DLB), avoidance of anticholinergic medications, and low-dose levodopa for parkinsonism. Prognosis is progressive.

\subsection{Frontotemporal Dementia}
\label{sec:geriatric:Frontotemporal Dementia}
Frontotemporal dementia (FTD) is a group of neurodegenerative disorders characterized by progressive dysfunction of the frontal and temporal lobes. Although typically presenting at a younger age (45--65 years), FTD can occur in older adults. The behavioral variant (bvFTD) presents with early and progressive changes in personality, social conduct, and executive function---disinhibition, apathy, loss of empathy, compulsive behaviors, and hyperorality. Language variants include non-fluent/agrammatic primary progressive aphasia (PPA) and semantic variant PPA. Approximately 30--50\% of FTD is familial, with mutations in MAPT (tau), GRN (progranulin), and C9orf72. There is no disease-modifying therapy. Management is symptomatic and supportive, focusing on behavior management (SSRIs for disinhibition), caregiver support, and addressing safety and legal concerns. Prognosis is progressive.

\subsection{Mild Cognitive Impairment}
\label{sec:Mild Cognitive Impairment}

Mild cognitive impairment (MCI) is an intermediate state between normal cognitive aging and dementia, characterized by objective cognitive decline (typically memory---amnestic MCI---or other domains) that is greater than expected for age and education but does not significantly impair functional independence. Prevalence is 10--20\% in adults aged 65 and older. The condition results from underlying neurodegenerative pathology, with the APOE epsilon-4 allele as a major genetic risk factor. Annual conversion rate to dementia is 5--15\%, with amnestic MCI conferring highest risk for Alzheimer disease. Clinical features include subjective cognitive complaints with objective deficits on testing. Proposed biomarkers include CSF amyloid-beta 42 (decreased), phosphorylated tau (increased), amyloid PET, and tau PET. Management includes cardiovascular risk factor control, physical exercise, cognitive stimulation, treatment of depression and sleep disorders, and avoidance of anticholinergic medications; no pharmacotherapy is approved for MCI. Prognosis is variable; annual conversion rate to dementia is 5--15\%.

\subsection{Normal Pressure Hydrocephalus}
\label{sec:Normal Pressure Hydrocephalus}

Normal pressure hydrocephalus (NPH) is a potentially reversible cause of dementia and gait impairment, characterized by the classic triad of gait apraxia (magnetic gait---feet appear glued to the floor), cognitive decline (subcortical pattern), and urinary incontinence (late symptom). Prevalence increases with age: up to 5--6\% of nursing home residents. Ventriculomegaly is present on brain imaging without significant cortical atrophy, and lumbar puncture opening pressure is normal ($\leq$20 cm HH$_2$O). The condition results from impaired CSF absorption. Clinical features include the triad of gait apraxia, cognitive decline, and urinary incontinence. Diagnosis is supported by clinical improvement after high-volume lumbar puncture (CSF tap test) and Evans index >0.3. Treatment is ventriculoperitoneal shunt placement, which improves gait in 70--80\% of appropriately selected patients. Prognosis is favorable with treatment; patient selection is critical for good outcomes.

\subsection{Parkinson Disease}
\label{sec:Parkinson Disease}
Parkinson disease (PD) is a progressive neurodegenerative disorder affecting 1--2\% of adults over 65. In the geriatric population, PD is associated with higher rates of postural instability, falls, dementia, dysphagia, and medication complications (motor fluctuations, dyskinesias, psychosis). Geriatric-specific considerations include management of orthostatic hypotension (common with levodopa and age-related autonomic dysfunction), careful use of dopamine agonists (increased risk of impulse control disorders, hallucinations, and sedation in older adults), and switching to levodopa monotherapy when possible. Deep brain stimulation (DBS) is effective but requires careful patient selection in older adults. Falls are a major concern, and physical therapy (LSVT BIG, balance training) is essential. Prognosis is progressive.

\subsection{Vascular Dementia}
\label{sec:Vascular Dementia}
Vascular dementia (VaD) is a progressive neurodegenerative disorder and the second most common dementia, resulting from cerebrovascular disease---large vessel infarcts (multi-infarct dementia), small vessel disease (subcortical ischemic vascular dementia), strategic single infarcts, cerebral amyloid angiopathy, or CADASIL. In the geriatric patient, hypertension, diabetes, atrial fibrillation, and smoking are the main modifiable risk factors. Clinical features include stepwise deterioration, patchy cognitive deficits (executive function disproportionately affected), focal neurological signs, and often a history of stroke or TIA. Vascular dementia frequently coexists with Alzheimer disease (mixed dementia). Brain MRI shows white matter hyperintensities, lacunar infarcts, and atrophy. Management focuses on vascular risk factor control (strict blood pressure control, antiplatelet therapy, statins, glycemic control) and symptomatic treatment of cognitive symptoms. Prognosis is progressive; geriatric patients require careful blood pressure management to avoid orthostatic hypotension and falls.

%-------------------------------------------------------------
\section{Falls and Mobility Disorders}
%-------------------------------------------------------------

%-------------------------------------------------------------

%-------------------------------------------------------------

%-------------------------------------------------------------

%-------------------------------------------------------------

%-------------------------------------------------------------

%-------------------------------------------------------------

\label{sec:Falls and Mobility Disorders}
%-------------------------------------------------------------%-------------------------------------------------------------

\subsection{Falls in the Elderly}
\label{sec:Falls in the Elderly}

Falls are the leading cause of fatal and non-fatal injuries in adults aged 65 and older. Approximately 30\% of community-dwelling older adults fall each year, increasing to 50\% in those over 80. Risk factors are classified as intrinsic (muscle weakness, gait/balance impairment, visual impairment, cognitive impairment, orthostatic hypotension, polypharmacy) and extrinsic (environmental hazards). Clinical features include fall-related injuries, fear of falling, and functional decline. The STEADI toolkit recommends annual fall screening. Prevention is multifactorial: strength and balance exercises (Tai Chi, Otago Exercise Programme), medication review and deprescribing, vitamin D supplementation, vision optimization, home hazard modification, and hip protectors for high-risk individuals. Prognosis is variable; multifactorial intervention reduces falls by 20--30\%.

\subsection{Fear of Falling}
\label{sec:Fear of Falling}

Fear of falling is a common psychological consequence occurring in 20--60\% of older adults after a fall, and can occur even without a history of falls. It leads to activity restriction, social isolation, and deconditioning (the ``post-fall syndrome''), paradoxically increasing the risk of future falls and functional decline. Assessment uses the Falls Efficacy Scale---International (FES-I) or the Activities-specific Balance Confidence (ABC) Scale. Management includes cognitive-behavioral therapy, graded exposure to feared activities, balance training (Tai Chi, yoga), and addressing the underlying physical deficits that perpetuate fear. Prognosis is favorable with appropriate intervention.

\subsection{Gait Disorders}
\label{sec:Gait Disorders}

Gait disorders affect 15\% of adults aged 60 and over and 80\% of those over 85. Common patterns include:
\begin{itemize}[noitemsep]
  \item \textbf{Parkinsonian gait:} Shuffling, reduced step length, decreased arm swing, festination (accelerating steps), freezing, and difficulty initiating gait. Associated with Parkinson disease and atypical parkinsonian disorders.
  \item \textbf{Hemiparetic gait:} Unilateral leg weakness with circumduction (leg swings in a semicircle), foot drop, and arm held in flexion. Seen in stroke and other unilateral brain lesions.
  \item \textbf{Ataxic gait:} Wide-based, unsteady, staggering, with inconsistent step width. Cerebellar etiology (stroke, tumor, alcohol-related, multiple system atrophy). Also seen in sensory ataxia (impaired proprioception from peripheral neuropathy or dorsal column disease) where vision partially compensates (positive Romberg sign).
  \item \textbf{Neuropathic gait} (steppage gait): Foot drop with high-stepping to clear the toe, often slapping of the foot. Caused by peripheral neuropathy (diabetic, alcoholic, Charcot-Marie-Tooth), L5 radiculopathy, or common peroneal nerve palsy.
  \item \textbf{Senile gait} (cautious gait): Slightly wide-based, shorter steps, reduced arm swing, mild stooped posture, and general caution. It may represent normal aging or a prodrome of neurological disease.
\end{itemize}
Evaluation includes history (onset, falls, neurological symptoms, medications), physical examination (neurological, musculoskeletal, vascular), and gait observation. Management addresses underlying cause, assistive devices (canes, walkers), physical therapy (gait training, balance exercises), and fall prevention strategies.

%-------------------------------------------------------------
\section{Frailty and Geriatric Syndromes}
%-------------------------------------------------------------

%-------------------------------------------------------------

%-------------------------------------------------------------

%-------------------------------------------------------------

%-------------------------------------------------------------

%-------------------------------------------------------------

%-------------------------------------------------------------

\label{sec:Frailty and Geriatric Syndromes}
%-------------------------------------------------------------%-------------------------------------------------------------

\subsection{Failure to Thrive}
\label{sec:Failure to Thrive}

Failure to thrive (FTT) in older adults is a geriatric syndrome of unintentional weight loss (>5\% in 6--12 months), decreased appetite, poor nutrition, and functional decline, often accompanied by social withdrawal and depression. It is multifactorial, resulting from medical (malignancy, dementia, depression, infection, hyperthyroidism, malabsorption, dysphagia), social (isolation, poverty, caregiver loss), and psychological (bereavement, delirium) causes. Clinical features include weight loss, poor appetite, functional decline, and social withdrawal. Diagnosis is clinical with identification of underlying causes. Treatment targets underlying causes, nutritional supplementation, appetite stimulants (megestrol acetate, mirtazapine), and social support. Prognosis depends on the underlying cause.

\subsection{Frailty}
\label{sec:Frailty}

Frailty is a geriatric syndrome of increased vulnerability to stressors resulting from age-related decline in physiological reserve across multiple organ systems. It affects a significant proportion of older adults. Two major conceptual frameworks guide frailty assessment. The \textbf{Fried Frailty Phenotype} (2001) defines frailty by the presence of three or more of five criteria: unintentional weight loss (>10 lb in one year), self-reported exhaustion, weakness (low grip strength), slow gait speed, and low physical activity; pre-frailty is defined as one or two criteria. The \textbf{Frailty Index} (Rockwood and Mitnitski) accumulates deficits across multiple domains as a ratio of deficits present to total deficits assessed (typically 30--70 items). The condition results from age-related decline in physiologic reserve, with sarcopenia as a core biological substrate. Clinical features include unintentional weight loss, exhaustion, weakness, slow gait speed, and low physical activity. Diagnosis uses the Fried Frailty Phenotype (three or more of five criteria) or the Frailty Index (>0.25 indicates frailty). Management includes resistance and aerobic exercise, protein-calorie supplementation, vitamin D, reduction of polypharmacy, and comprehensive geriatric assessment (CGA). Prognosis is variable; frailty is a dynamic process that can improve or worsen over time.

\subsection{Malnutrition in the Elderly}
\label{sec:Malnutrition in the Elderly}

Malnutrition is a geriatric syndrome affecting 15--60\% of hospitalized older adults and 30--50\% of nursing home residents. The condition results from physiologic changes (reduced olfaction and taste, delayed gastric emptying), social factors (isolation, poverty, inability to shop/cook), dentition problems, dysphagia, polypharmacy, depression, and dementia. Clinical features include weight loss, poor nutritional intake, and functional decline. Diagnosis uses the Mini Nutritional Assessment (MNA) as a validated screening tool. Management includes dietary counseling, oral nutritional supplements, texture-modified foods for dysphagia, treatment of depression, dental care, and in severe cases, enteral nutrition. Prognosis is variable; consequences include impaired immune function, delayed wound healing, muscle wasting, falls, and increased mortality.

\subsection{Pressure Ulcers}
\label{sec:Pressure Ulcers}

Pressure ulcers (bedsores, decubitus ulcers) are localized injury to skin and underlying tissue caused by prolonged pressure, shear, or friction, most commonly over bony prominences (sacrum, heels, greater trochanters, ischial tuberosities). Prevalence is 5--15\% in hospitalized patients and up to 25\% in long-term care. Risk factors include immobility, malnutrition, incontinence, sensory loss, advanced age, and comorbidities such as diabetes and peripheral vascular disease. Clinical features range from non-blanchable erythema (Stage 1) to full-thickness tissue loss with exposed bone/tendon/muscle (Stage 4), with stages defined by the National Pressure Injury Advisory Panel (NPIAP). Diagnosis is clinical with staging. Prevention is paramount: risk assessment (Braden Scale), pressure-redistribution surfaces, regular repositioning every 2 hours, nutritional support, and skin care. Treatment includes wound debridement, moisture management, infection control, negative-pressure wound therapy, and surgical flap closure for advanced cases. Prognosis is variable; healing depends on the stage and underlying patient condition.

\subsection{Sarcopenia}
\label{sec:Sarcopenia}

Sarcopenia is a geriatric syndrome defined as the progressive, generalized loss of skeletal muscle mass, strength, and function with aging, with prevalence of 5--13\% in adults aged 60--70 years and up to 50\% in those over 80. The condition results from neural denervation of motor units, reduced anabolic signaling (growth hormone, IGF-1, testosterone), increased catabolic cytokines (IL-6, TNF-alpha), mitochondrial dysfunction, and inadequate protein intake. Clinical features include low muscle mass, weakness, and functional decline. The European Working Group on Sarcopenia in Older People (EWGSOP2) defines sarcopenia by low muscle strength (grip strength <27 kg men, <16 kg women) plus low muscle quantity/quality and confirmed by low physical performance (gait speed $\leq$0.8 m/s). Treatment is resistance exercise (2--3 times weekly, progressive overload) plus adequate protein intake ($\ge$1.2 g/kg/day) and vitamin D (800 IU/day); testosterone replacement may be considered in hypogonadal men. Prognosis is variable; early intervention can improve outcomes.

%-------------------------------------------------------------
\section{Geriatric Palliative and End-of-Life Care}
%-------------------------------------------------------------

%-------------------------------------------------------------

%-------------------------------------------------------------

%-------------------------------------------------------------

%-------------------------------------------------------------

%-------------------------------------------------------------

%-------------------------------------------------------------

\label{sec:Geriatric Palliative and End-of-Life Care}
%-------------------------------------------------------------%-------------------------------------------------------------

\subsection{Advance Care Planning}
\label{sec:Advance Care Planning}

Advance care planning (ACP) is the ongoing process of discussing values, goals, and preferences for future medical care with patients and their surrogates before the patient loses decision-making capacity. Components include designation of a healthcare proxy, advance directives, do-not-resuscitate (DNR) orders, do-not-intubate (DNI) orders, and Physician Orders for Life-Sustaining Treatment (POLST forms). ACP reduces unwanted hospitalizations, improves goal-concordant care, and reduces distress for families and surrogates. In older adults with dementia, ACP should be initiated early in the disease course while the patient retains decisional capacity. Serious illness conversations should address prognosis, expected functional decline, goals of care, and specific treatment preferences.

\subsection{Ethical Issues}
\label{sec:Ethical Issues}

Geriatric care raises complex ethical challenges:
\begin{itemize}[noitemsep]
  \item \textbf{Decision-making capacity:} Capacity is task-specific and can be assessed using the Aid to Capacity Evaluation (ACE) tool. A patient with capacity can understand relevant information, appreciate the situation and its consequences, reason about treatment options, and communicate a choice. Incapacity requires surrogate decision-making based on substituted judgment or best interests.
  \item \textbf{Informed consent:} Must be obtained from patients with capacity; for those lacking capacity, consent is obtained from the legally authorized representative.
  \item \textbf{Truth-telling:} Geriatric patients should be informed of their diagnosis and prognosis in a sensitive, individualized manner. Withholding information (therapeutic privilege) is rarely justified.
  \item \textbf{Euthanasia and assisted dying:} Medical assistance in dying (MAiD) is legal in several jurisdictions. In Canada, MAiD is available to adults with a grievous and irremediable medical condition causing enduring suffering, with a reasonably foreseeable natural death. Advance requests for MAiD in dementia require careful ethical and legal consideration.
  \item \textbf{Elder abuse and neglect:} Affects 10--15\% of older adults, including physical, emotional, financial abuse, and neglect. Clinicians have an ethical and legal duty to report suspected abuse to adult protective services.
\end{itemize}

\subsection{Hospice Care}
\label{sec:Hospice Care}

Hospice care is a model of interdisciplinary, palliative care for patients with terminal illness and a life expectancy of 6 months or less (certified by two physicians), focused on comfort and quality of life rather than curative treatment. Hospice is provided in the home, nursing home, or hospice facility. Eligibility for common geriatric conditions: dementia (Functional Assessment Staging---FAST stage 7), heart failure, COPD, and cancer. The interdisciplinary team includes physicians, nurses, social workers, chaplains, volunteers, and bereavement counselors. Despite its benefits, hospice is underused in older adults; median length of stay is 18 days, and 30\% of patients die within 7 days of enrollment. Earlier referral improves outcomes and family satisfaction.

\subsection{Pain Management in the Elderly}
\label{sec:Pain Management in the Elderly}

Chronic pain affects 50--75\% of older adults, most commonly from osteoarthritis, degenerative spine disease, neuropathic conditions, and musculoskeletal disorders. Pain is often under-recognized and under-treated in older adults, especially those with cognitive impairment (use of PAINAD scale or Abbey Pain Scale for non-verbal patients). The WHO analgesic ladder guides treatment, but geriatric-specific cautions apply:
\begin{itemize}[noitemsep]
  \item \textbf{Acetaminophen (paracetamol):} First-line for musculoskeletal pain; maximum 3000 mg/day to avoid hepatotoxicity.
  \item \textbf{NSAIDs:} Avoid chronic use in older adults due to gastrointestinal, renal, and cardiovascular risks; use lowest dose, shortest duration if unavoidable, with PPI for GI protection.
  \item \textbf{Topical analgesics:} Lidocaine patches, topical NSAIDs (diclofenac), capsaicin---preferred for localized pain with minimal systemic side effects.
  \item \textbf{Opioids:} Reserved for moderate-to-severe pain not responding to non-opioid therapies; start low, go slow; monitor for constipation, sedation, respiratory depression, and falls.
  \item \textbf{Gabapentinoids (gabapentin, pregabalin):} Effective for neuropathic pain; renally dosed; associated with sedation, dizziness, falls.
  \item \textbf{Adjuvants:} SNRIs (duloxetine) and TCAs (nortriptyline preferred over amitriptyline for lower anticholinergic burden) for neuropathic pain and fibromyalgia.
\end{itemize}

%-------------------------------------------------------------
\section{Incontinence}
%-------------------------------------------------------------

%-------------------------------------------------------------

%-------------------------------------------------------------

Geriatric incontinence encompasses urinary and fecal involuntary leakage, significantly affecting quality of life, dignity, and independence in older adults. Etiologies involve complex interactions between detrusor muscle dysfunction, pelvic floor weakness, neurological impairment, and environmental factors. Multimodal management combines behavioral retraining, pelvic muscle therapy, targeted pharmacological agents, and supportive environmental modifications.

%-------------------------------------------------------------

%-------------------------------------------------------------

%-------------------------------------------------------------

%-------------------------------------------------------------

\label{sec:Incontinence}
%-------------------------------------------------------------%-------------------------------------------------------------

\subsection{Fecal Incontinence}
\label{sec:Fecal Incontinence}

Fecal incontinence (FI)---involuntary loss of solid or liquid stool---affects 10--20\% of community-dwelling older adults and up to 50\% of nursing home residents. Causes include impaired anal sphincter function (obstetric trauma, surgical injury, age-related degeneration), reduced rectal sensation (diabetes, spinal cord injury, fecal impaction with overflow), diarrhea, and cognitive/mobility impairment. Diagnosis uses the St. Mark's (Vaizey) score to quantify severity. Management includes dietary modification (fiber supplementation for formed stool, avoidance of dietary triggers), scheduled toilet routines, pelvic floor rehabilitation, antidiarrheal medications, anal plugs, and surgical options for severe refractory cases. Addressing constipation and fecal impaction in institutionalized patients is critical.

\subsection{Urinary Incontinence}
\label{sec:geriatric:Urinary Incontinence}
Urinary incontinence (UI) is a geriatric syndrome affecting 30--60\% of older adults, increasing with age and more common in women. Types include stress (leakage with coughing, sneezing, laughing---urethral hypermobility or intrinsic sphincter deficiency), urge (overactive bladder---leakage with sudden strong urge to void, due to detrusor overactivity), overflow (incomplete bladder emptying with continuous dribbling---often from outlet obstruction or detrusor underactivity), functional (inability to reach toilet due to mobility impairment, cognitive deficit, or environmental barriers), and mixed. In the geriatric population, transient causes must be excluded (DIAPPERS mnemonic: Delirium, Infection, Atrophic urethritis/vaginitis, Pharmaceuticals, Psychological, Excess urine output, Restricted mobility, Stool impaction). Diagnosis includes history, bladder diary, post-void residual by ultrasound, urinalysis, and cough stress test. Management includes behavioral therapy (pelvic floor exercises, bladder training, scheduled voiding, fluid management), medications (antimuscarinics or beta-3 agonists for overactive bladder), absorbent products, pessaries, and surgical options for refractory cases. Prognosis is variable.

%-------------------------------------------------------------
\section{Osteoporosis and Bone Health}
%-------------------------------------------------------------

%-------------------------------------------------------------

%-------------------------------------------------------------

%-------------------------------------------------------------

%-------------------------------------------------------------

%-------------------------------------------------------------

%-------------------------------------------------------------

\label{sec:Osteoporosis and Bone Health}
%-------------------------------------------------------------%-------------------------------------------------------------

\subsection{Age-Related Bone Loss}
\label{sec:Age-Related Bone Loss}
Osteoporosis is a systemic skeletal disorder characterized by low bone mass and microarchitectural deterioration of bone tissue, leading to increased fracture risk. It affects 10--20\% of women over 50 and 4--5\% of men; prevalence increases steeply with age. The condition results from an imbalance between bone resorption and bone formation, driven by estrogen deficiency, secondary hyperparathyroidism, decreased mechanical loading, and reduced anabolic hormone levels. Diagnosis is by dual-energy X-ray absorptiometry (DXA): a T-score of $\leq$-2.5 defines osteoporosis; -1.0 to -2.5 defines osteopenia. The FRAX tool estimates 10-year fracture probability. Treatment includes calcium and vitamin D supplementation, weight-bearing exercise, fall prevention, bisphosphonates (first-line), denosumab, teriparatide, and romosozumab. Prognosis is favorable with treatment.

\subsection{Hip Fracture}
\label{sec:Hip Fracture}

Hip fracture is one of the most devastating injuries in older adults, with a 1-year mortality of 20--30\% and up to 50\% of survivors losing independent mobility. The incidence increases exponentially with age, doubling every 5--6 years after age 65. Osteoporotic hip fractures typically result from a fall from standing height or less. Anatomically, hip fractures are classified as femoral neck (intracapsular) or intertrochanteric (extracapsular). Management requires urgent orthopedic consultation, surgical repair within 24--48 hours (associated with reduced mortality and complications), and comprehensive perioperative care including pain management, delirium prevention, thromboprophylaxis, infection prophylaxis, early mobilization, nutritional support, and osteoporosis treatment. Post-operative rehabilitation with physical therapy is essential for functional recovery. Prognosis is guarded; 1-year mortality is 20--30\%.

%-------------------------------------------------------------
\section{Polypharmacy and Adverse Drug Reactions}
%-------------------------------------------------------------

%-------------------------------------------------------------

%-------------------------------------------------------------

%-------------------------------------------------------------

%-------------------------------------------------------------

%-------------------------------------------------------------

%-------------------------------------------------------------

\label{sec:Polypharmacy and Adverse Drug Reactions}
%-------------------------------------------------------------%-------------------------------------------------------------

\subsection{Beers Criteria / STOPP/START Criteria}
\label{sec:Beers Criteria / STOPP/START Criteria}

Two major tools guide prescribing in older adults. The \textbf{Beers Criteria} (American Geriatrics Society, updated 2023) lists potentially inappropriate medications (PIMs) to avoid or use with caution in older adults, including anticholinergics, benzodiazepines, Z-drugs, NSAIDs, sulfonylureas, and PPIs. The \textbf{STOPP/START criteria} are a European alternative that identifies PIMs (STOPP) and potential prescribing omissions (START---e.g., missing statin in diabetes with cardiovascular disease, missing bisphosphonate after osteoporotic fracture). Both criteria reduce inappropriate prescribing and ADRs when integrated into clinical practice.

\subsection{Drug Interactions in the Elderly}
\label{sec:Drug Interactions in the Elderly}

Age-related physiological changes alter pharmacokinetics (absorption, distribution---increased body fat, decreased total body water, decreased albumin---metabolism, and elimination---reduced renal function) and pharmacodynamics (increased sensitivity to CNS-active drugs, anticholinergics, anticoagulants). Common drug-drug interactions in older adults include warfarin---NSAIDs (increased bleeding risk), warfarin---amiodarone (increased INR), ACE inhibitors---spironolactone (hyperkalemia), SSRIs---NSAIDs (GI bleeding), and opioids---benzodiazepines (respiratory depression). Management includes maintaining an up-to-date medication list, using a single pharmacy and prescriber when possible, prioritizing drug-drug interaction checking, and avoiding high-risk combinations.

\subsection{Medication Non-Adherence}
\label{sec:Medication Non-Adherence}

Medication non-adherence affects 40--60\% of older adults, contributing to poor disease control, increased hospitalizations, and mortality. Barriers include cognitive impairment (forgetting doses), complex regimens (multiple pills, multiple daily doses), cost, transportation difficulties, poor provider-patient communication, lack of understanding of medication purpose, and side effects. Interventions include simplifying regimens (once-daily dosing, combination pills), medication organizers (blister packs, pillboxes), caregiver involvement, clear written instructions, technology-based reminders, and addressing polypharmacy through deprescribing.

\subsection{Polypharmacy}
\label{sec:Polypharmacy}

Polypharmacy---commonly defined as the concurrent use of five or more medications---affects 40--70\% of older adults. It is associated with increased risk of adverse drug reactions (ADRs), drug-drug interactions, medication non-adherence, falls, cognitive impairment, hospitalization, and mortality. Factors driving polypharmacy include multimorbidity, multiple prescribers, lack of medication reconciliation, and prescribing cascades (where ADRs are misinterpreted as new medical conditions and treated with additional medications). Management requires regular medication review, deprescribing (systematic process of discontinuing potentially inappropriate medications), and the use of explicit criteria to guide prescribing decisions.

%-------------------------------------------------------------
\section{Sensory Impairment in Aging}
%-------------------------------------------------------------

%-------------------------------------------------------------

%-------------------------------------------------------------

%-------------------------------------------------------------

%-------------------------------------------------------------

%-------------------------------------------------------------

%-------------------------------------------------------------

\label{sec:Sensory Impairment in Aging}
%-------------------------------------------------------------%-------------------------------------------------------------

\subsection{Age-Related Macular Degeneration}
\label{sec:geriatric:Age-Related Macular Degeneration}
Age-related macular degeneration (AMD) is the leading cause of irreversible central vision loss in adults over 65. Dry (nonexudative) AMD accounts for 85--90\% of cases, characterized by drusen, retinal pigment epithelium atrophy, and photoreceptor degeneration. Wet (exudative/neovascular) AMD is less common but more aggressive, driven by choroidal neovascularization. In older adults, AMD severely impacts quality of life---loss of reading, face recognition, and driving independence---and is associated with depression and falls. Management for dry AMD includes the AREDS2 formulation (vitamins C and E, lutein/zeaxanthin, zinc) for intermediate disease; wet AMD is treated with intravitreal anti-VEGF injections. Prognosis is guarded; low-vision rehabilitation is essential for all stages.

\subsection{Presbycusis}
\label{sec:Presbycusis}
Presbycusis is age-related hearing loss affecting two-thirds of adults over 70. It is a sensorineural hearing loss, typically bilateral, symmetrical, and high-frequency (above 2000 Hz), resulting from degeneration of cochlear hair cells, stria vascularis, and spiral ganglion neurons. In older adults, untreated hearing loss is associated with cognitive decline (increased dementia risk), social isolation, depression, and falls. Management includes hearing aids (rehabilitative benefit is well established), cochlear implants for severe-to-profound loss, communication strategies, and treating cerumen impaction (common, easily reversible cause of hearing loss in older adults). Prognosis is favorable with amplification.

\subsection{Presbyopia and Cataracts}
\label{sec:Presbyopia and Cataracts}
Presbyopia is the age-related loss of accommodation resulting from lens hardening, affecting virtually all adults over 45, corrected with reading glasses, bifocals, or progressive lenses. Cataracts---opacification of the crystalline lens---affect 50\% of adults over 75 and are the leading cause of reversible vision loss worldwide, classified by location: nuclear sclerotic (most common, progressive myopic shift), cortical (spoke-like opacities), and posterior subcapsular (associated with corticosteroid use, diabetes, and radiation). Surgical removal with intraocular lens (IOL) implantation is definitive treatment, with phacoemulsification being the standard technique, yielding >95\% improvement in visual acuity. Prognosis is excellent.

